\documentclass[../../main.tex]{subfiles}

\graphicspath{{../../image/}{../../figure/}}

\begin{document}

\section{$BV$ 函数理论概述}\label{section:ch01-BV函数理论概述}

$BV$ 函数空间为带激波阵面的平衡律方程组解提供了自然框架。事实上，这类函数的一个突出特征是：其间断点聚集在余维 $1$ 的流形上。$BV$ 函数理论的完整论述可在\cref{section:ch01-注记} 所引文献中找到，此处仅列出与我们目的相关的性质，不给证明。

\begin{definition}\label{definition:ch01-局部BV}
$\mathbb{R}^{k}$ 开子集 $\mathscr{X}$ 上的标量函数 $v$ 称为在 $\mathscr{X}$ 上\emph{局部有界变差}，若 $v\in L^{1}_{\mathrm{loc}}(\mathscr{X})$ 且 $\operatorname{grad}v$ 是 $\mathscr{X}$ 上一个（取值于 $\mathbb{R}^{k}$ 的）Radon 测度 $\mathscr{M}$，即对任一检验函数 $\Psi\in C_{0}^{\infty}(\mathscr{X};\mathbb{M}^{1\times k})$，
\begin{align}
-\int_{\mathscr{X}}v\,\operatorname{div}\Psi(X)\,\mathrm{d}X=\int_{\mathscr{X}}\Psi(X)\,\mathrm{d}\mathscr{M}(X).\label{align:ch01-BV-def}
\end{align}
当 $v\in L^{1}(\mathscr{X})$ 且 $\mathscr{M}$ 有限时，$v$ 是 $\mathscr{X}$ 上的\emph{有界变差函数}，其\emph{全变差}为
\begin{align}
TV_{\mathscr{X}}v=|\mathscr{M}|(\mathscr{X})=\sup_{|\Psi(X)|=1}\int_{\mathscr{X}}v(X)\operatorname{div}\Psi(X)\,\mathrm{d}X.\label{align:ch01-total-variation}
\end{align}
$\mathscr{X}$ 上有界变差与局部有界变差函数的集合分别记为 $BV(\mathscr{X})$ 与 $BV_{\mathrm{loc}}(\mathscr{X})$。
\end{definition}

显然，Sobolev 空间 $W^{1,1}(\mathscr{X})$（即 $L^{1}(\mathscr{X})$ 中导数亦在 $L^{1}(\mathscr{X})$ 的函数）包含于 $BV(\mathscr{X})$；而 $W^{1,1}_{\mathrm{loc}}(\mathscr{X})$ 包含于 $BV_{\mathrm{loc}}(\mathscr{X})$。

下述命题给出检验一个给定函数是否具有有界变差的有用准则：

\begin{theorem}\label{theorem:ch01-差商准则}
设 $\{E_{\alpha}:\alpha=1,\cdots,k\}$ 表示 $\mathbb{R}^{k}$ 的标准正交基。若 $v\in BV_{\mathrm{loc}}(\mathscr{X})$，则对任一开有界集 $\mathscr{Y}$（$\overline{\mathscr{Y}}\subset\mathscr{X}$），
\begin{align}
\limsup_{h\downarrow 0}\frac{1}{h}\int_{\mathscr{Y}}|v(X+hE_{\alpha})-v(X)|\,\mathrm{d}X=|\mathscr{M}_{\alpha}|(\mathscr{Y}),\qquad \alpha=1,\cdots,k,\label{align:ch01-difference-quotient}
\end{align}
其中 $\mathscr{M}_{\alpha}$ 为 $\mathscr{M}$ 的相应分量。反之，若 $v\in L^{1}_{\mathrm{loc}}(\mathscr{X})$ 且对上述每个 $\mathscr{Y}$，\cref{align:ch01-difference-quotient} 左端有限，则 $v\in BV_{\mathrm{loc}}(\mathscr{X})$。
\end{theorem}

作为推论，上述命题给出如下紧性结果：

\begin{theorem}\label{theorem:ch01-BV紧性}
$BV_{\mathrm{loc}}(\mathscr{X})$ 中任一序列 $\{v_{\ell}\}$，若对每个开有界 $\mathscr{Y}\subset\mathscr{X}$，$\|v_{\ell}\|_{L^{1}(\mathscr{Y})}$ 与 $TV_{\mathscr{Y}}v_{\ell}$ 一致有界，则存在子序列在 $L^{1}_{\mathrm{loc}}(\mathscr{X})$ 中、且在 $\mathscr{X}$ 上几乎处处收敛于某 $v\in BV_{\mathrm{loc}}(\mathscr{X})$，满足 $TV_{\mathscr{Y}}v\leqslant\liminf_{\ell\to\infty}TV_{\mathscr{Y}}v_{\ell}$。
\end{theorem}

有界变差函数具有精细的几何结构，如下述定理所述：

\begin{theorem}\label{theorem:ch01-BV-三分}
任一 $v\in BV_{\mathrm{loc}}(\mathscr{X})$ 的定义域可分解为三个两两不相交子集 $\mathscr{C}$、$\mathscr{J}$、$\mathscr{I}$ 之并，具有下列性质：
\begin{enumerate}
\item\label{item:ch01-BV-三分-c} $\mathscr{C}$ 是 $v$ 的\emph{近似连续点}集，即对每个 $\bar{X}\in\mathscr{C}$，伴随 $v_{0}\in\mathbb{R}$ 使得
\begin{align}
\lim_{r\downarrow 0}\frac{1}{r^{k}}\int_{\mathscr{B}_{r}(\bar{X})}|v(X)-v_{0}|\,\mathrm{d}X=0.\label{align:ch01-approx-continuity}
\end{align}
\item\label{item:ch01-BV-三分-j} $\mathscr{J}$ 是 $v$ 的\emph{近似跳跃间断点}集，即对每个 $\bar{X}\in\mathscr{J}$，伴随 $N\in\mathbb{S}^{k-1}$ 与不同的 $v_{-},v_{+}\in\mathbb{R}$，使得
\begin{align}
\lim_{r\downarrow 0}\frac{1}{r}\int_{\mathscr{B}_{r}^{\pm}(\bar{X})}|v(X)-v_{\pm}|\,\mathrm{d}X=0,\label{align:ch01-approx-jump}
\end{align}
其中 $\mathscr{B}_{r}^{\pm}(\bar{X})$ 表示半球 $\mathscr{B}_{r}(\bar{X})\cap\{X:(X-\bar{X})\cdot N\gtrless 0\}$。此外，$\mathscr{J}$ 是可数可求长的，即它本质上被一族嵌入 $\mathbb{R}^{k}$ 的 $C^{1}$ 类 $(k-1)$ 维流形 $\{\mathscr{F}_{i}\}$ 的可数并所覆盖：$\mathscr{H}^{k-1}(\mathscr{J}\setminus\bigcup_{i}\mathscr{F}_{i})=0$。而且当 $\bar{X}\in\mathscr{J}\cap\mathscr{F}_{i}$ 时，$N$ 在 $\bar{X}$ 处是 $\mathscr{F}_{i}$ 的法向。
\item\label{item:ch01-BV-三分-i} $\mathscr{I}$ 是 $v$ 的\emph{不规则点}集；其 $(k-1)$ 维 Hausdorff 测度为零：$\mathscr{H}^{k-1}(\mathscr{I})=0$。
\end{enumerate}
\end{theorem}

至此，$BV$ 函数的身份在一 $(k$ 维 Lebesgue$)$ 零测集上修改其值后不受影响，即 $BV_{\mathrm{loc}}(\mathscr{X})$ 实际上是函数等价类空间，仅在差一零测集的意义下确定。然而，当处理这些函数更精细的行为时，宜为每个等价类指定一个规范代表，其值至多差一 $(k-1)$ 维 Hausdorff 零测集。具体做法如下。

设 $g$ 为 $\mathbb{R}$ 上的连续函数，$v\in BV_{\mathrm{loc}}(\mathscr{X})$。参照\cref{theorem:ch01-BV-三分} 的记号，$g$ 与 $v$ 的\emph{规范化复合} $\widetilde{g\circ v}$ 定义为
\begin{align}
\widetilde{g\circ v}(X)=
\begin{cases}
g(v_{0}), & \text{若 }X\in\mathscr{C},\\[1mm]
\displaystyle\int_{0}^{1}g(\tau v_{-}+(1-\tau)v_{+})\,\mathrm{d}\tau, & \text{若 }X\in\mathscr{J},
\end{cases}\label{align:ch01-normalized-composition}
\end{align}
而在不规则点集 $\mathscr{I}$（其 $(k-1)$ 维 Hausdorff 测度为零）上任意取值。特别地，可对 $v$ 自身规范化：
\begin{align}
\tilde{v}(X)=
\begin{cases}
v_{0}, & \text{若 }X\in\mathscr{C},\\[1mm]
\frac{1}{2}(v_{-}+v_{+}), & \text{若 }X\in\mathscr{J}.
\end{cases}\label{align:ch01-normalized-v}
\end{align}
于是 $\mathscr{C}$ 的每一点都成为 Lebesgue 点。

上述规范化的适当性由经典链式法则的下述推广所表明：

\begin{theorem}\label{theorem:ch01-BV链式法则}
设 $g$ 在 $\mathbb{R}$ 上连续可微，导数为 $\mathsf{D}g$，$v\in BV_{\mathrm{loc}}(\mathscr{X})\cap L^{\infty}(\mathscr{X})$。则 $g\circ v\in BV_{\mathrm{loc}}(\mathscr{X})\cap L^{\infty}(\mathscr{X})$。规范化函数 $\widetilde{\mathsf{D}g\circ v}$ 关于测度 $\mathscr{M}=\operatorname{grad}v$ 局部可积，且
\begin{align}
\operatorname{grad}(g\circ v)=\widetilde{\mathsf{D}g\circ v}\,\operatorname{grad}v\label{align:ch01-chain-rule}
\end{align}
在如下意义下：对任一检验函数 $\Psi\in C_{0}^{\infty}(\mathscr{X};\mathbb{M}^{1\times k})$，
\begin{align}
-\int_{\mathscr{X}}g(v(X))\operatorname{div}\Psi(X)\,\mathrm{d}X=\int_{\mathscr{X}}\widetilde{\mathsf{D}g\circ v}(X)\Psi(X)\,\mathrm{d}\mathscr{M}(X).\label{align:ch01-chain-rule-weak}
\end{align}
\end{theorem}

下面考察 $\mathbb{R}^{k}$ 中与 $BV$ 函数理论密切相关的一类集合的重要几何性质。

\begin{definition}\label{definition:ch01-有限周长}
$\mathbb{R}^{k}$ 的子集 $\mathscr{D}$ 称为具有（局部）有限周长，若其指示函数 $\chi_{\mathscr{D}}$ 在 $\mathbb{R}^{k}$ 上具有（局部）有界变差。
\end{definition}

将\cref{theorem:ch01-BV-三分} 应用于具有局部有限周长的集合 $\mathscr{D}$ 的指示函数 $\chi_{\mathscr{D}}$。显然，$\chi_{\mathscr{D}}$ 的近似连续点集 $\mathscr{C}$ 是 $\mathscr{D}$ 与 $\mathbb{R}^{k}\setminus\mathscr{D}$ 的密度点集之并。$\mathscr{C}$ 的补集，即 $\mathbb{R}^{k}$ 中既非 $\mathscr{D}$ 密度点亦非 $\mathbb{R}^{k}\setminus\mathscr{D}$ 密度点的点 $X$ 全体，构成 $\mathscr{D}$ 的\emph{测度论边界} $\partial\mathscr{D}$。可以证明，$\mathscr{D}$ 具有有限周长当且仅当 $\mathscr{H}^{k-1}(\partial\mathscr{D})<\infty$，其周长可由 $TV_{\mathbb{R}^{k}}\chi_{\mathscr{D}}$ 或 $\mathscr{H}^{k-1}(\partial\mathscr{D})$ 度量。$\chi_{\mathscr{D}}$ 的近似跳跃间断点集称为 $\mathscr{D}$ 的\emph{约化边界}，记为 $\partial^{*}\mathscr{D}$。由\cref{theorem:ch01-BV-三分}，$\partial^{*}\mathscr{D}\subset\partial\mathscr{D}$、$\mathscr{H}^{k-1}(\partial\mathscr{D}\setminus\partial^{*}\mathscr{D})=0$，且 $\partial^{*}\mathscr{D}$ 被可数个 $C^{1}$ 类 $(k-1)$ 维流形之并覆盖。此外，与 $\partial^{*}\mathscr{D}$ 上每一点 $X$ 相联系的向量 $N\in\mathbb{S}^{k-1}$，可自然解释为 $\mathscr{D}$ 在 $X$ 处的\emph{测论外法向}。具 Lipschitz 边界的集合有有限周长。事实上，平衡律的整个理论可以这样重构：把真域取为\cref{section:ch01-平衡律的表述} 中所假定的不必 Lipschitz、而仅具有限周长的集合。

\begin{definition}\label{definition:ch01-迹}
设 $\mathscr{D}$ 具有有限周长，$v\in BV_{\mathrm{loc}}(\mathbb{R}^{k})$。在 $\mathscr{D}$ 的约化边界 $\partial^{*}\mathscr{D}$ 上外法向为 $N$ 的点 $\bar{X}$ 处，$v$ 具有\emph{内迹}与\emph{外迹} $v_{-}$、$v_{+}$，若
\begin{align}
\lim_{r\downarrow 0}\frac{1}{r^{k}}\int_{\mathscr{B}_{r}^{\pm}}|v(X)-v_{\pm}|\,\mathrm{d}X=0.\label{align:ch01-trace-def}
\end{align}
\end{definition}

可以证明，迹 $v_{\pm}$ 对 $\partial^{*}\mathscr{D}$ 上几乎所有（关于 $\mathscr{H}^{k-1}$）点有定义，并在 $\partial^{*}\mathscr{D}$ 上局部可积。此外，成立如下形式的 Gauss--Green 定理：

\begin{theorem}\label{theorem:ch01-Gauss-Green}
设 $v\in BV(\mathbb{R}^{k})$，故 $\mathscr{M}=\operatorname{grad}v$ 为有限测度。考虑任一有界、具有限周长的集合 $\mathscr{D}$，其密度点集为 $\mathscr{D}^{*}$、约化边界为 $\partial^{*}\mathscr{D}$。则
\begin{align}
\mathscr{M}(\mathscr{D}^{*})=\int_{\partial^{*}\mathscr{D}}v_{+}N\,\mathrm{d}\mathscr{H}^{k-1}.\label{align:ch01-GG-interior}
\end{align}
此外，对 $\partial\mathscr{D}$ 的任一 Borel 子集 $\mathscr{F}$，
\begin{align}
\mathscr{M}(\mathscr{F})=\int_{\mathscr{F}}(v_{-}-v_{+})N\,\mathrm{d}\mathscr{H}^{k-1}.\label{align:ch01-GG-jump}
\end{align}
\end{theorem}

特别地，$BV_{\mathrm{loc}}(\mathbb{R}^{k})$ 中任一 $v$ 的近似跳跃间断点集 $\mathscr{J}$ 可被可数个定向曲面之并覆盖，因而\cref{align:ch01-GG-jump} 对 $\mathscr{J}$ 的任一可测子集 $\mathscr{F}$ 成立。

对 $v\in BV(\mathscr{X})$，测度 $\mathscr{M}=\operatorname{grad}v$ 可分解为三个相互奇异测度之和：其\emph{连续部}，关于 $k$ 维 Lebesgue 测度绝对连续；其\emph{跳跃部}，集中在 $v$ 的近似跳跃间断点集 $\mathscr{J}$ 上；以及其\emph{Cantor 部}。特别地，$\mathscr{X}$ 的任一 Borel 子集上测度的 Cantor 部，当其 $(k-1)$ 维 Hausdorff 测度有限时消失。

\begin{definition}\label{definition:ch01-SBV}
$v\in BV(\mathscr{X})$ 称为\emph{特殊有界变差函数}，记作 $v\in SBV(\mathscr{X})$，若测度 $\operatorname{grad}v$ 的 Cantor 部消失。
\end{definition}

可以证明，$SBV(\mathscr{X})$ 是 $BV(\mathscr{X})$ 的真子空间，并真包含 $W^{1,1}(\mathscr{X})$。

当 $k=1$ 时，$BV$ 函数理论与经典的有界变差函数理论密切相关。设 $v$ 为区间 $(a,b)\subset(-\infty,\infty)$（有界或无界）上的 $BV$ 函数。令 $\tilde{v}$ 为 $v$ 的规范化形式。则
\begin{align}
TV_{(a,b)}v=\sup\sum_{j=1}^{\ell-1}|\tilde{v}(x_{j+1})-\tilde{v}(x_{j})|,\label{align:ch01-TV-1d}
\end{align}
其中上确界取遍所有（有限）网 $a<x_{1}<x_{2}<\cdots<x_{\ell}<b$。此外，（经典）单侧极限 $\tilde{v}(x\pm)$ 在每一点 $x\in(a,b)$ 存在，且除至多可数个点外二者均等于 $\tilde{v}(x)$。当 $k=1$ 时，紧性\cref{theorem:ch01-BV紧性} 化为经典的 Helly 定理。

任一 $v\in SBV(a,b)$ 是一个绝对连续函数与一个跳跃函数之和。相应地，测度 $\operatorname{grad}v$ 是 $v$ 的点态导数 $v'$（在 $(a,b)$ 上几乎处处存在）与（至多可数个）位于 $v$ 跳跃间断点、权重为跃度的加权 Dirac 质量之和。

向量值函数 $U$ 在 $\mathscr{X}$ 上（局部）有界变差，若其每个分量在 $\mathscr{X}$ 上都具有（局部）有界变差；其全变差 $TV_{\mathscr{X}}U$ 为各分量全变差之和。上述关于标量值函数的所有讨论，特别是\cref{theorem:ch01-差商准则,theorem:ch01-BV紧性,theorem:ch01-BV-三分,theorem:ch01-BV链式法则,theorem:ch01-Gauss-Green} 的断言，均可立即推广到（并将在下文用于）向量值有界变差函数。

\end{document}
